\section{Contacts at many depths and completion of the proof}
\label{sec:contact-lower}

We continue to assume coexistence in the coordinate direction and adapt the
cut-averaging argument of \cite[Proposition~7.1]{OpenAI2026} to scripted cuts
and positive-side proximity labels.  The first-contact estimate from
Proposition~\ref{prop:first-contact} forces contacts in most groups
of consecutive dyadic depth intervals.  We will first count these intervals and
then use the finite mean width of the common blocks to convert depths into block
indices.  The resulting lower bound contradicts
Proposition~\ref{prop:contact-upper}.

\begin{proposition}[Contact lower bound]\label{prop:contact-lower}
Suppose that $P_0(A_{e_1})>0$ and $P_0(A_{-e_1})>0$, and let $m(J)$ be the
contact statistic of the opposed tapes defined in \eqref{eq:contact-count}.
There are constants $c>0$ and an integer $c_1\ge0$ such that, for all sufficiently
large integers $J$,
\begin{equation}\label{eq:contact-lower}
 m(J+c_1)\ge \frac{cJ}{1+\log\log J}.
\end{equation}
\end{proposition}

\begin{proof}
All constants in this proof may depend on the fixed environment law and the
marked construction, but not on $J$ or the depth parameters.  Recall that the
depth of a contact is the depth of its positive departure.  For an integer
$l\ge0$, call $l$ a \emph{contact depth label} if there is a contact at an integer
depth in $[2^l,2^{l+1})$.

\paragraph{Conditioning on a distant cut.}
Fix a large integer $J$, and set
\begin{equation}\label{eq:lower-scales}
 N=2^{2J},\qquad
 G=\left\lceil 4\log_2\log(2+J)\right\rceil+3.
\end{equation}
Let $\mathcal C_N$ be the event that $N$ is a width renewal of the positive
tape, and write
\[
 \mathbf Q_N=\mathbf P(\,\cdot\mid\mathcal C_N).
\]
By \eqref{eq:bridge-renewal},
\begin{equation}\label{eq:lower-conditioning}
 \mathbf P(\mathcal C_N)=\bar u_N^+\ge\xi p_+>0.
\end{equation}
Under $\mathbf Q_N$, follow the placed positive words from bottom to top by
reversing the order of the slab list through this renewal, without reversing
any word, and translate its bottom endpoint to the origin.  The product formula
\eqref{eq:bridge-law} is invariant under this list
reversal.  Thus the resulting chronological path $Y$ has bridge law
$\mathcal B_N^+$.  If $E_{\rm lat}=\pi Y_{T_N}$, the same translation puts the
negative tape at $(N-1,E_{\rm lat})$.  Write $\beta$ for this translated path.
Its displacement path from its starting site has law $\widehat P_-$ and is
independent of the positive list.  This is
exactly the pair of paths in Proposition~\ref{prop:first-contact}.  In these
coordinates the depth of a positive departure at $y$ is $N-y_1$.

On $\mathcal C_N$, at most $N$ positive slabs are used before width $N$, since
each width is at least $M\ge1$.  The norm of $E_{\rm lat}$ is bounded by the sum
of their radii, and hence by the sum of the first $N$ radii of the unconditioned
positive tape.  Proposition~\ref{prop:radius-moment},
\eqref{eq:lower-conditioning}, and Markov's inequality give
\begin{equation}\label{eq:lower-endpoint-tail}
 \mathbf Q_N\{\abs{E_{\rm lat}}>N^2\}\le \frac{C}{N}.
\end{equation}

\paragraph{An interval without contacts forces a late first contact.}
For $G<l\le J$, let $U_l'$ be the event that none of
$l-G,l-G+1,\ldots,l$ is a contact depth label.  We will show that
\begin{equation}\label{eq:uncovered-labels}
 \sum_{l=G+1}^J\mathbf Q_N(U_l')=o(J).
\end{equation}
Consider the possible positive cut depths
\[
 \mathcal I_l=[2^l,\tfrac32\,2^l]\cap\ZZ.
\]
For $s\in\mathcal I_l$, put
\begin{equation}\label{eq:lower-contact-parameters}
 k=N-s,\qquad r=N-2^{l-G}.
\end{equation}
For all sufficiently large $J$, these integers satisfy
$M\le k<r<N$ and $r-k\ge R$, uniformly over the indicated $l,s$.

Suppose first that $s$ is an actual cut depth of the positive tape.  In the
reversed list, its boundary is at height $k$.  Every preceding word remains
strictly below its terminal height until its final arrival, and that arrival
ends a record script.  Consequently the bridge first reaches $k$ along a
prefix realizing $\mathcal D_k^+$: it first hits $k-M$ and then follows the
prescribed $M$-step script.  All the subsequent words remain at or above $k$.
These assertions hold for the concatenated bridge, because a word's initial
boundary is a strict record of the preceding words and no script can straddle
a boundary at which another script ends.

The final positive departure before $T_N$ is $(N-1,E_{\rm lat})$, the initial
site of $\beta$.  There is therefore a first contact along $Y$ between $T_k$
and $T_N$.  Its depth is at most $s$.  On $U_l'$, there are no contacts at any
depth in $[2^{l-G},2^{l+1})$; since $s<2^{l+1}$, this first contact must have
depth less than $2^{l-G}$.  Its height is in particular at least $r$.  If also
$\abs{E_{\rm lat}}\le N^2$, all the requirements of
\eqref{eq:first-contact-event} are now satisfied.

For each \emph{deterministic} $s\in\mathcal I_l$, let $\mathcal A_{l,s}$ denote
the event \eqref{eq:first-contact-event} with the parameters
\eqref{eq:lower-contact-parameters}.  This event includes the prefix condition
$\mathcal D_k^+$, but we do not additionally require $s$ to be a cut of the
tape.  Thus Proposition~\ref{prop:first-contact} applies directly under
$\mathbf Q_N$, and gives
\begin{equation}\label{eq:lower-deterministic-cut}
 \mathbf Q_N(\mathcal A_{l,s})
 \le C\log(2N)\,\phi(\Delta_{l,s}),\qquad
 \Delta_{l,s}=u^+_{s-2^{l-G}}-u_s^+,
\end{equation}
where
\[
 \phi(x)=x\log^2(\mathrm e/x)\quad(0<x\le1),\qquad \phi(0)=0.
\]
Here and below $\mathrm e=\exp(1)$.  Lemma~\ref{lem:width-tail} implies
\begin{equation}\label{eq:lower-delta-tail}
 \Delta_{l,s}\le C2^{-G}a_l,\qquad
 a_l=2^l F_+(2^{l-1}),\qquad
 A:=\sum_{l\ge1}a_l<\infty,
\end{equation}
because $s-2^{l-G}\ge2^{l-1}$.
The factor $2^{-G}$ is the ratio between the remaining contact depth
$2^{l-G}$ and the cut depth scale $2^l$.

\paragraph{Averaging over the cut locations.}
Let $C_+(t)$ count the positive width renewals in $(0,t]$.  The positive widths
are iid with finite mean by Propositions~\ref{prop:cut-law}
and~\ref{prop:mean-width}; their strong law gives
$C_+(t)=t/E_{\nu_+}L+o(t)$ almost surely.  It follows that the number of positive
cuts in $\mathcal I_l$ is
\[
 \frac{2^{l-1}}{E_{\nu_+}L}+o(2^l)
 \quad\text{almost surely as }l\longrightarrow\infty.
\]
Fix $0<c_0<1/(2E_{\nu_+}L)$, and let $\mathcal T_J$ be the event that every
$\mathcal I_l$ with $G<l\le J$ contains at least $c_0 2^l$ positive cuts.  On
almost every tape the displayed asymptotic supplies this bound for all
sufficiently large $l$.  Since $G\to\infty$,
\[
 \mathbf P(\mathcal T_J^c)\longrightarrow0,
 \qquad
 \mathbf Q_N(\mathcal T_J^c)
 \le\frac{\mathbf P(\mathcal T_J^c)}{\xi p_+}\longrightarrow0.
\]
In particular, no convergence rate in the strong law is needed for the
changing conditional law $\mathbf Q_N$.

Every cut counted in $\mathcal I_l$ lies before depth $N$.  On $U_l'$,
$\mathcal T_J$, and $\{\abs{E_{\rm lat}}\le N^2\}$, each of these cuts supplies
one of the events $\mathcal A_{l,s}$ by the preceding argument.  We obtain the
pointwise inequality
\begin{equation}\label{eq:cut-averaging}
 \one_{U_l'}\one_{\mathcal T_J}\one_{\{\abs{E_{\rm lat}}\le N^2\}}
 \le\frac{1}{c_0 2^l}\sum_{s\in\mathcal I_l}\one_{\mathcal A_{l,s}}.
\end{equation}
The right side involves only the deterministic-parameter probabilities already
bounded in \eqref{eq:lower-deterministic-cut}.

We next sum those bounds using only $\sum_l a_l<\infty$.  For large $J$,
$C2^{-G}a_l\le\mathrm e^{-1}$ uniformly in $l$, and $\phi$ is increasing on
$[0,\mathrm e^{-1}]$.  Hence \eqref{eq:lower-delta-tail} yields
\begin{equation}\label{eq:lower-phi}
 \phi(\Delta_{l,s})
 \le C2^{-G}a_l\bigl((G+1)^2+\log_+^2(1/a_l)\bigr),
\end{equation}
where $\log_+ x=\max\{0,\log x\}$ and the right side is interpreted as zero
when $a_l=0$.  Moreover,
\begin{equation}\label{eq:lower-slow-tail}
 \sum_{l=1}^J a_l\log_+^2(1/a_l)
 \le C(A+J^{-1})\log^2 J.
\end{equation}
Indeed, terms with $a_l\ge J^{-2}$ contribute at most $4A\log^2 J$.
For the remaining terms, the function $x\log^2(1/x)$ is increasing on
$[0,J^{-2}]$ for large $J$, so each is at most $4J^{-2}\log^2 J$.

Taking expectations in \eqref{eq:cut-averaging}, using
$|\mathcal I_l|\le2^l$, and applying
\eqref{eq:lower-endpoint-tail}, \eqref{eq:lower-deterministic-cut},
\eqref{eq:lower-phi}, and \eqref{eq:lower-slow-tail}, we obtain
\begin{align}
 \sum_{l=G+1}^J\mathbf Q_N(U_l')
 &\le J\mathbf Q_N(\mathcal T_J^c)+\frac{CJ}{N}\notag\\
 &\quad+C\log(2N)\,2^{-G}
       \bigl(A(G+1)^2+(A+J^{-1})\log^2 J\bigr)
 =o(J).\label{eq:lower-uncovered-bound}
\end{align}
For the final equality, the first term is $o(J)$, while
\eqref{eq:lower-scales} gives $\log(2N)=O(J)$,
$2^{-G}=O((\log J)^{-4})$, and $G=O(\log\log J)$.
The last term is therefore $O(J/(\log J)^2)$, and $CJ/N=o(J)$.  This proves
\eqref{eq:uncovered-labels}, even when the summable sequence $(a_l)$ decays
arbitrarily slowly.

There are $J-G$ tested labels $l$.  By \eqref{eq:uncovered-labels} and Markov's
inequality, with $\mathbf Q_N$-probability tending to one, at least $J/2$ of
these tests have $U_l'$ false.  Each such test has a contact depth label in
$\{l-G,\ldots,l\}$, and any one contact label can account for at most $G+1$
tests.  Thus, if $\mathcal E_J$ is the event that at least $J/(2(G+1))$ of
$1,\ldots,J$ are contact depth labels, then
\begin{equation}\label{eq:lower-unconditioning}
 \mathbf Q_N(\mathcal E_J)\longrightarrow1,
 \qquad
 \liminf_{J\to\infty}\mathbf P(\mathcal E_J)\ge\xi p_+>0.
\end{equation}
The second assertion follows by intersecting $\mathcal E_J$ with
$\mathcal C_N$ and using \eqref{eq:lower-conditioning}.  We have therefore
obtained the needed number of depth labels under the original tape law.

\paragraph{From depths to common-block indices.}
By Proposition~\ref{prop:common-blocks}, $\mu=\mathbf E K_1$ is finite and positive,
and $W_i'/i\to\mu$ almost surely.  For an integer depth $z\ge1$, let $i(z)$ be
its common-block index, so that
\[
 W_{i(z)-1}'<z\le W_{i(z)}'.
\]
The strong law implies, almost surely for all sufficiently large integer $z$,
\begin{equation}\label{eq:depth-index-comparison}
 \frac{z}{2\mu}\le i(z)<\frac{2z}{\mu}+1.
\end{equation}
Consequently there is a deterministic integer $c_1\ge0$ such that, almost
surely for all sufficiently large $l$, every depth $2^l\le z<2^{l+1}$ has
an index label $j$ satisfying
\[
 2^j<i(z)\le2^{j+1},\qquad |j-l|\le c_1.
\]
The opposite choices of open endpoints for the depth and index intervals
affect this comparison by at most one label and are included in $c_1$.
Let $\mathcal V_J$ be the event that the comparison holds for every
$l\ge\lceil\sqrt J\rceil$.  Its probability tends to one.  By
\eqref{eq:lower-unconditioning}, the intersection $\mathcal E_J\cap\mathcal V_J$
therefore has probability bounded below by a positive constant for all large
$J$.

On this intersection discard the depth labels below $\lceil\sqrt J\rceil$.
This loses only $O(\sqrt J)=o(J/(G+1))$ labels.  Each remaining contact depth
label gives a contact in a block-index interval labeled in
$\{1,\ldots,J+c_1\}$, and each index label can receive at most $2c_1+1$
distinct depth labels.  There are consequently at least $cJ/(G+1)$ such
index labels on an event of probability bounded below.  Taking expectations
in the definition \eqref{eq:contact-count} proves
\[
 m(J+c_1)\ge\frac{cJ}{G+1}.
\]
Since $G+1=O(1+\log\log J)$, this is \eqref{eq:contact-lower}.
\end{proof}

\begin{proof}[Proof of Theorem~\ref{thm:main}]
If both coordinate escape events $A_{e_1}$ and $A_{-e_1}$ had positive
probability, Propositions~\ref{prop:contact-upper} and
\ref{prop:contact-lower} would give, for all large $J$,
\[
 \frac{cJ}{1+\log\log J}
 \le m(J+c_1)
 \le \frac{C(J+c_1)}{\log(J+c_1)},
\]
which is impossible as $J\to\infty$.  The same argument applies after any
permutation of the coordinate axes, so coexistence is impossible in every
coordinate direction.

Now fix any nonzero real direction $\ell$.  If $0<P_0(A_\ell)<1$,
Proposition~\ref{prop:sign-union} gives positive probability to both
$A_\ell$ and $A_{-\ell}$.  Proposition~\ref{prop:coordinate-reduction} then
gives a coordinate direction with positive probability for both escape
signs, contradicting what we have just proved.  Thus $P_0(A_\ell)\in\{0,1\}$.
\end{proof}
